% =====================================================================
%  IEL04 - Circuitos Eléctricos I
%  Semana 1: El condensador: estructura, capacitancia, tipos e identificación de componentes comerciales
%  Institución Universitaria de Barranquilla (IUB)
%  Generada por src/presentaciones.py (diseño IUB 5). Compilador: pdfLaTeX (dos pasadas)
% =====================================================================
\documentclass[aspectratio=169,11pt,xcolor={table}]{beamer}

% ---------------------------------------------------------------------
%  Codificación, idioma y tipografía
% ---------------------------------------------------------------------
\usepackage[utf8]{inputenc}
\usepackage[T1]{fontenc}
\usepackage[spanish,es-nodecimaldot,es-noshorthands]{babel}
\usepackage{avant}                       % Sans geométrica (emula Avenir Next)
\renewcommand{\familydefault}{\sfdefault}
\renewcommand{\ttdefault}{pcr}           % Monoespaciada para código
\usepackage{textcomp}

% ---------------------------------------------------------------------
%  Paquetes
% ---------------------------------------------------------------------
\usepackage{amsmath,amssymb,mathtools}
\usepackage{siunitx}
\sisetup{output-decimal-marker={.}, per-mode=symbol}
\usepackage{booktabs,tabularx,multirow,array}
\usepackage{adjustbox}
\usepackage{tikz}
\usetikzlibrary{arrows.meta,positioning,calc,shapes.geometric,shapes.misc,
  decorations.pathreplacing,fit,backgrounds,babel}
\usepackage[siunitx,RPvoltages]{circuitikz}
\usepackage{pgfplots}
\pgfplotsset{compat=1.18}
\usepackage[most]{tcolorbox}
\usepackage{listings}

% ---------------------------------------------------------------------
%  Paleta institucional IUB (Manual de Identidad Visual 2025)
% ---------------------------------------------------------------------
\definecolor{iubwhite}{HTML}{FFFFFF}
\definecolor{iubnavy}{HTML}{121023}
\definecolor{iubyellow}{HTML}{FFDF2D}
\definecolor{iubblue}{HTML}{00ADE7}
\definecolor{iuborange}{HTML}{E39037}
\definecolor{iubgreen}{HTML}{3B9C6D}

% ---------------------------------------------------------------------
%  Configuración de Beamer
% ---------------------------------------------------------------------
\setbeamertemplate{navigation symbols}{}
\setbeamersize{text margin left=0.6cm, text margin right=0.6cm}
\setbeamercolor{normal text}{fg=iubnavy, bg=iubwhite}
\setbeamercolor{background canvas}{bg=iubwhite}
\setbeamercolor{structure}{fg=iubnavy}
\setbeamercolor{alerted text}{fg=iuborange}
\setbeamertemplate{itemize item}{\textcolor{iubblue}{\small$\blacktriangleright$}}
\setbeamertemplate{itemize subitem}{\textcolor{iuborange}{\scriptsize$\bullet$}}
\setbeamertemplate{enumerate item}{\textcolor{iubblue}{\bfseries\insertenumlabel.}}
\setbeamertemplate{frametitle}{}         % el título vive en el headline

% Parches: el headline se pre-mide antes del primer frame
\makeatletter
\providecommand{\insertframetitle}{}
\providecommand{\insertframesubtitle}{}
\makeatother

% ---------- Encabezado ----------
\setbeamertemplate{headline}{%
  \begin{tikzpicture}
    \useasboundingbox (0,0) rectangle (\paperwidth,1.05cm);
    \fill[iubnavy] (0,0) rectangle (\paperwidth,1.05cm);
    \fill[iubyellow] (0,0) rectangle (\paperwidth,0.05cm);
    \node[anchor=west, text=iubyellow, font=\large\bfseries]
      at (0.6cm,0.55cm) {\strut\insertframetitle};
    \node[anchor=east, text=iubyellow, font=\footnotesize\bfseries]
      at ([xshift=-0.6cm]\paperwidth,0.55cm) {IUB};
  \end{tikzpicture}}

% ---------- Pie de página ----------
\setbeamertemplate{footline}{%
  \begin{tikzpicture}
    \useasboundingbox (0,0) rectangle (\paperwidth,0.55cm);
    \fill[iubnavy] (0,0) rectangle (\paperwidth,0.55cm);
    \node[anchor=west, text=iubyellow, font=\tiny]
      at (0.6cm,0.275cm) {IEL04 \textperiodcentered{} Circuitos Eléctricos I};
    \node[anchor=center, text=iubyellow, font=\tiny]
      at (0.66\paperwidth,0.275cm) {Ing. Sergio Martinez Castilla};
    \node[anchor=east, text=iubyellow, font=\tiny\bfseries]
      at ([xshift=-0.6cm]\paperwidth,0.275cm)
      {Semana 1 \quad \insertframenumber\,/\,\inserttotalframenumber};
  \end{tikzpicture}}

% ---------------------------------------------------------------------
%  Cajas tcolorbox (texto blanco sobre la paleta complementaria)
%  \begin{caja}[opciones]{color}{título} ... \end{caja}
%  \begin{cajaplana}[opciones]{color} ... \end{cajaplana}
%  \begin{cardB|cardO|cardG|cardN}[opciones]{título} ... \end{cardX}
% ---------------------------------------------------------------------
\tcbset{
  iubbase/.style={enhanced, arc=1.5mm, boxrule=0pt, coltext=white, coltitle=white,
    fonttitle=\bfseries\footnotesize, fontupper=\footnotesize,
    left=1.8mm, right=1.8mm, top=1mm, bottom=1.2mm,
    toptitle=0.6mm, bottomtitle=0.6mm, halign=left,
    before upper={\hyphenpenalty=5000\setbeamercolor{itemize item}{fg=white}%
                  \setbeamercolor{itemize/enumerate body}{fg=white}%
                  \setbeamercolor{itemize/enumerate subbody}{fg=white}%
                  \setbeamercolor{itemize subitem}{fg=white}%
                  \setbeamercolor{enumerate item}{fg=white}%
                  \setbeamertemplate{itemize item}{\textcolor{white}{\small$\blacktriangleright$}}}}
}
\newtcolorbox{caja}[3][]{iubbase, colback=#2, colframe=#2,
  colbacktitle=#2!78!iubnavy, title={#3}, #1}
\newtcolorbox{cajaplana}[2][]{iubbase, colback=#2, colframe=#2, #1}
\newtcolorbox{cardB}[2][]{iubbase, colback=iubblue,   colframe=iubblue,
  colbacktitle=iubblue!78!iubnavy,   title={#2}, #1}
\newtcolorbox{cardO}[2][]{iubbase, colback=iuborange, colframe=iuborange,
  colbacktitle=iuborange!78!iubnavy, title={#2}, #1}
\newtcolorbox{cardG}[2][]{iubbase, colback=iubgreen,  colframe=iubgreen,
  colbacktitle=iubgreen!78!iubnavy,  title={#2}, #1}
\newtcolorbox{cardN}[2][]{iubbase, colback=iubnavy,   colframe=iubnavy,
  colbacktitle=iubnavy, coltitle=iubyellow, title={#2}, #1}

% Celda de encabezado de tabla (texto amarillo sobre \rowcolor{iubnavy}) y código en línea
\newcommand{\hd}[1]{\textcolor{iubyellow}{\bfseries #1}}
\newcommand{\thc}[1]{\textcolor{iubyellow}{\textbf{#1}}}
\newcommand{\cod}[1]{\texttt{\textbf{#1}}}

% ---------------------------------------------------------------------
%  Código sobre fondo oscuro:
%  \begin{codebox}[listing options app={language=Python,firstnumber=1}]{título}
% ---------------------------------------------------------------------
\lstdefinestyle{iubcode}{
  basicstyle=\ttfamily\fontsize{7}{8.4}\selectfont\color{white},
  keywordstyle=\color{iubblue}\bfseries,
  commentstyle=\color{iubgreen!55!white}\itshape,
  stringstyle=\color{iuborange!80!white},
  numbers=left, numberstyle=\tiny\color{iubyellow}, numbersep=6pt,
  showstringspaces=false, breaklines=true, tabsize=4,
  columns=fullflexible, keepspaces=true, upquote=true}
\newtcblisting{codebox}[2][]{enhanced, listing only, colback=iubnavy,
  colframe=iubnavy, colbacktitle=iubnavy!85!white, coltitle=iubyellow,
  fonttitle=\bfseries\scriptsize, title={#2}, arc=1.5mm, boxrule=0pt,
  left=6mm, right=1mm, top=0.5mm, bottom=0.5mm, toptitle=0.5mm, bottomtitle=0.5mm,
  listing options={style=iubcode}, #1}

% ---------------------------------------------------------------------
%  Estilos TikZ
% ---------------------------------------------------------------------
\tikzset{
  bloque/.style={rectangle, rounded corners=2mm, fill=#1, text=white,
    font=\scriptsize\bfseries, align=center, minimum height=1.1cm,
    text width=1.8cm, inner sep=1.2mm},
  flecha/.style={-{Stealth[length=2.2mm]}, line width=1pt, draw=iubnavy},
  arr/.style={flecha},
  term/.style={rectangle, draw=#1, line width=1.2pt, fill=white,
    text=iubnavy, font=\tiny\bfseries, minimum width=1.05cm,
    minimum height=0.5cm, inner sep=1pt},
  nodo/.style={rectangle, rounded corners=1mm, fill=#1, text=white,
    font=\scriptsize\bfseries, minimum size=0.75cm, align=center},
  cable/.style={line width=1.4pt, draw=#1},
  box/.style={rectangle, rounded corners=1.5mm, draw=none, text=white,
    font=\scriptsize\bfseries, align=center, minimum height=0.8cm, inner sep=3pt},
  bB/.style={box, fill=iubblue}, bO/.style={box, fill=iuborange},
  bG/.style={box, fill=iubgreen}, bN/.style={box, fill=iubnavy},
  dec/.style={diamond, aspect=1.9, fill=iuborange, text=white,
    font=\scriptsize\bfseries, align=center, inner sep=1pt},
  tag/.style={fill=#1, text=white, font=\scriptsize\bfseries, rounded corners=1mm,
    minimum width=0.7cm, anchor=north west},
}

% ---------------------------------------------------------------------
%  Diapositiva separadora de bloque
%  #1 número  #2 título  #3 duración  #4 color
% ---------------------------------------------------------------------
\newcommand{\bloque}[4]{%
\section{#2}
\begin{frame}[plain]
\begin{tikzpicture}[remember picture, overlay]
  \fill[#4] (current page.north west) rectangle
        ([xshift=5.2cm]current page.south west);
  \node[anchor=north west, text=white, font=\bfseries\large]
    at ([xshift=0.8cm,yshift=-2.2cm]current page.north west) {BLOQUE};
  \node[anchor=north west, text=white, font=\bfseries\fontsize{72}{72}\selectfont]
    at ([xshift=0.65cm,yshift=-2.9cm]current page.north west) {#1};
  \node[anchor=north west, text=iubnavy, font=\bfseries\LARGE,
        text width=9.4cm, align=left]
    at ([xshift=6.2cm,yshift=-2.8cm]current page.north west)
    {\hyphenpenalty=10000\exhyphenpenalty=10000 #2};
  \node[anchor=north west, fill=iubnavy, text=iubyellow, rounded corners=1.5mm,
        font=\small\bfseries, inner sep=2mm]
    at ([xshift=6.2cm,yshift=-6.0cm]current page.north west) {Duración: #3};
  \fill[iubnavy] ([xshift=5.2cm]current page.south west) rectangle
        ([yshift=0.35cm]current page.south east);
  \fill[iubyellow] ([xshift=5.2cm,yshift=0.35cm]current page.south west)
        rectangle ([yshift=0.42cm]current page.south east);
  \node[anchor=north east, text=iubnavy, font=\small\bfseries]
    at ([xshift=-0.6cm,yshift=-0.5cm]current page.north east) {IUB · IEL04};
\end{tikzpicture}
\end{frame}}

% =====================================================================
\begin{document}
% =====================================================================

% ---------------------------------------------------------------------
%  PORTADA
% ---------------------------------------------------------------------
\begin{frame}[plain]
\begin{tikzpicture}[remember picture, overlay]
  % Panel institucional
  \fill[iubnavy] (current page.north west) rectangle
        ([xshift=5.6cm]current page.south west);
  \node[anchor=north west, text=iubyellow, font=\bfseries\fontsize{54}{54}\selectfont]
    at ([xshift=0.7cm,yshift=-1.0cm]current page.north west) {IUB};
  \node[anchor=north west, text=white, font=\footnotesize, text width=4.3cm, align=left]
    at ([xshift=0.75cm,yshift=-3.0cm]current page.north west)
    {Institución Universitaria de Barranquilla};
  \draw[iubyellow, line width=1.2pt]
    ([xshift=0.75cm,yshift=-4.25cm]current page.north west) --
    ([xshift=2.6cm,yshift=-4.25cm]current page.north west);
  \node[anchor=north west, text=iubyellow, font=\scriptsize\bfseries,
        text width=4.3cm, align=left]
    at ([xshift=0.75cm,yshift=-4.5cm]current page.north west)
    {Tecnología en Gestión de Sistemas Eléctricos};
  \node[anchor=south west, text=white, font=\scriptsize, text width=4.3cm, align=left]
    at ([xshift=0.75cm,yshift=0.9cm]current page.south west)
    {Ing. Sergio Martinez Castilla\\[1pt]\textcolor{iubyellow}{\tiny smartinezc@unibarranquilla.edu.co}};
  % Bloque de título
  \node[anchor=north west, fill=iubblue, text=white, rounded corners=1.5mm,
        font=\scriptsize\bfseries, inner sep=2mm]
    at ([xshift=6.4cm,yshift=-1.0cm]current page.north west)
    {IEL04 \textperiodcentered{} Semana 1 de 14 \textperiodcentered{} Corte evaluativo 1};
  \node[anchor=north west, text=iubnavy, font=\small, text width=9.0cm, align=left] (a)
    at ([xshift=6.4cm,yshift=-1.9cm]current page.north west)
    {Circuitos Eléctricos I};
  \node[anchor=north west, text=iubnavy, font=\bfseries\fontsize{15}{18}\selectfont,
        text width=9.0cm, align=left] (t)
    at ([yshift=-0.35cm]a.south west)
    {\hyphenpenalty=10000 El condensador: estructura, capacitancia, tipos e identificación de componentes comerciales};
  % Franjas de la paleta complementaria
  \fill[iubblue]   ([xshift=6.4cm,yshift=1.1cm]current page.south west)
    rectangle ++(1.6cm,0.18cm);
  \fill[iuborange] ([xshift=8.1cm,yshift=1.1cm]current page.south west)
    rectangle ++(1.6cm,0.18cm);
  \fill[iubgreen]  ([xshift=9.8cm,yshift=1.1cm]current page.south west)
    rectangle ++(1.6cm,0.18cm);
  \node[anchor=north west, text=iubnavy, font=\scriptsize]
    at ([xshift=6.4cm,yshift=0.95cm]current page.south west)
    {Sesión teórico-práctica \textperiodcentered{} 240 min};
\end{tikzpicture}
\end{frame}

% ---------------------------------------------------------------------
\begin{frame}{Punto de partida de la sesión}
\begin{columns}[T,onlytextwidth]
  \column{0.34\textwidth}
\begin{caja}[height=5.9cm]{iubblue}{Continuidad}
\scriptsize \begin{itemize}\setlength{\itemsep}{2pt}
  \item \textbf{Semana 1: Condensadores}
  \item Semana 2: Bobinas o inductores
\end{itemize}
\end{caja}
  \column{0.34\textwidth}
\begin{caja}[height=5.9cm]{iuborange}{Prerrequisitos}
\scriptsize \begin{itemize}\setlength{\itemsep}{2pt}
  \item Definir carga eléctrica, diferencia de potencial
  \item Aplicar la ley de Ohm a un resistor y calcular potencia en cd
  \item Usar notación científica y de ingeniería con prefijos métricos
  \item Leer el código de colores y la tolerancia de un resistor
\end{itemize}
\end{caja}
  \column{0.29\textwidth}
\begin{caja}[height=5.9cm]{iubgreen}{Competencia}
\scriptsize \textbf{UC2.} Coordinar actividades asociadas a sistemas eléctricos utilizando fuentes convencionales y no convencionales de energía.
\end{caja}
\end{columns}
\end{frame}

% ---------------------------------------------------------------------
\begin{frame}{Resultados de aprendizaje}
\begin{columns}[c,onlytextwidth]
  \column{0.45\textwidth}
\begin{caja}{iubnavy}{IEL04-RA1}
\footnotesize Calcular un circuito capacitivo, inductivo y LC, sea en serie o paralelo.
\end{caja}
\vspace{1mm}
\begin{cajaplana}{iubblue}
\scriptsize \textbf{CE1.} Identificar los tipos de capacitores e inductores.\\[2pt]
\textbf{CE2.} Realizar mediciones en un circuito capacitivo e inductivo.\\[2pt]
\textbf{CE3.} Realizar mediciones en un circuito LC.
\end{cajaplana}
  \column{0.52\textwidth}
\centering
\begin{adjustbox}{max width=\linewidth, max totalheight=6.1cm}
\begin{tikzpicture}
  \node[box, fill=iubblue, text width=4.9cm, align=left, font=\tiny, anchor=west] (r0) at (1.25,0.00) {Explicar la estructura interna de un condensador y el papel del dieléctrico en su capacitancia.};
  \node[tag=iubnavy, font=\tiny\bfseries, minimum width=1.15cm, anchor=east] at (1.1,0.00) {Comprender};
  \node[box, fill=iuborange, text width=4.9cm, align=left, font=\tiny, anchor=west] (r1) at (1.25,-1.45) {Calcular la capacitancia, la carga y la energía almacenada de un condensador a partir de su geometría o de su tensión de trabajo.};
  \node[tag=iubnavy, font=\tiny\bfseries, minimum width=1.15cm, anchor=east] at (1.1,-1.45) {Aplicar};
  \node[box, fill=iubgreen, text width=4.9cm, align=left, font=\tiny, anchor=west] (r2) at (1.25,-2.90) {Identificar el tipo, el valor nominal, la tolerancia y la tensión de trabajo de condensadores comerciales a partir de su marcado.};
  \node[tag=iubnavy, font=\tiny\bfseries, minimum width=1.15cm, anchor=east] at (1.1,-2.90) {Aplicar};
  \node[box, fill=iubblue, text width=4.9cm, align=left, font=\tiny, anchor=west] (r3) at (1.25,-4.35) {Medir la capacitancia de condensadores con un capacímetro y contrastarla con su intervalo de tolerancia, como base de las mediciones en circuitos capacitivos y LC.};
  \node[tag=iubnavy, font=\tiny\bfseries, minimum width=1.15cm, anchor=east] at (1.1,-4.35) {Evaluar};
  \draw[flecha, draw=iubnavy!40] (r0.south) -- (r1.north);
  \draw[flecha, draw=iubnavy!40] (r1.south) -- (r2.north);
  \draw[flecha, draw=iubnavy!40] (r2.south) -- (r3.north);
\end{tikzpicture}
\end{adjustbox}
\end{columns}
\end{frame}

% ---------------------------------------------------------------------
\begin{frame}{Ruta de la sesión \textperiodcentered{} 240 minutos}
\centering
\begin{tikzpicture}[x=1cm,y=1cm]
  \node[circle, fill=iubblue, text=white, font=\scriptsize\bfseries, minimum size=0.6cm, inner sep=0] at (0,0.00) {1};
  \node[anchor=west, font=\scriptsize, text width=7.6cm] at (0.45,0.00) {El condensador como elemento que almacena energía en el campo eléctrico y su lugar en los circuitos};
  \fill[iubblue] (8.3,-0.20) rectangle ++(2.04,0.4);
  \node[anchor=west, font=\scriptsize\bfseries] at (10.44,0.00) {20 min};
  \node[circle, fill=iuborange, text=white, font=\scriptsize\bfseries, minimum size=0.6cm, inner sep=0] at (0,-0.76) {2};
  \node[anchor=west, font=\scriptsize, text width=7.6cm] at (0.45,-0.76) {Placas, dieléctrico, C = Q/V y capacitancia de placas paralelas};
  \fill[iuborange] (8.3,-0.96) rectangle ++(3.58,0.4);
  \node[anchor=west, font=\scriptsize\bfseries] at (11.98,-0.76) {35 min};
  \node[circle, fill=iubgreen, text=white, font=\scriptsize\bfseries, minimum size=0.6cm, inner sep=0] at (0,-1.51) {3};
  \node[anchor=west, font=\scriptsize, text width=7.6cm] at (0.45,-1.51) {Separación de cargas, rigidez dieléctrica, tensión nominal y energía W = CV\texttwosuperior{}/2};
  \fill[iubgreen] (8.3,-1.71) rectangle ++(3.58,0.4);
  \node[anchor=west, font=\scriptsize\bfseries] at (11.98,-1.51) {35 min};
  \node[circle, fill=iubblue, text=white, font=\scriptsize\bfseries, minimum size=0.6cm, inner sep=0] at (0,-2.27) {4};
  \node[anchor=west, font=\scriptsize, text width=7.6cm] at (0.45,-2.27) {Cerámicos, de película, de mica, electrolíticos de aluminio y tantalio, variables};
  \fill[iubblue] (8.3,-2.47) rectangle ++(3.58,0.4);
  \node[anchor=west, font=\scriptsize\bfseries] at (11.98,-2.27) {35 min};
  \node[circle, fill=iuborange, text=white, font=\scriptsize\bfseries, minimum size=0.6cm, inner sep=0] at (0,-3.03) {5};
  \node[anchor=west, font=\scriptsize, text width=7.6cm] at (0.45,-3.03) {Lectura del valor, la tolerancia y la tensión en el cuerpo del condensador};
  \fill[iuborange] (8.3,-3.23) rectangle ++(3.07,0.4);
  \node[anchor=west, font=\scriptsize\bfseries] at (11.47,-3.03) {30 min};
  \node[circle, fill=iubgreen, text=white, font=\scriptsize\bfseries, minimum size=0.6cm, inner sep=0] at (0,-3.79) {6};
  \node[anchor=west, font=\scriptsize, text width=7.6cm] at (0.45,-3.79) {Faradio y sus submúltiplos, conversiones y símbolos normalizados};
  \fill[iubgreen] (8.3,-3.99) rectangle ++(2.56,0.4);
  \node[anchor=west, font=\scriptsize\bfseries] at (10.96,-3.79) {25 min};
  \node[circle, fill=iubblue, text=white, font=\scriptsize\bfseries, minimum size=0.6cm, inner sep=0] at (0,-4.54) {7};
  \node[anchor=west, font=\scriptsize, text width=7.6cm] at (0.45,-4.54) {Decodificar el marcado, calcular el intervalo de tolerancia y contrastarlo con la medición};
  \fill[iubblue] (8.3,-4.74) rectangle ++(4.60,0.4);
  \node[anchor=west, font=\scriptsize\bfseries] at (13.00,-4.54) {45 min};
  \node[circle, fill=iuborange, text=white, font=\scriptsize\bfseries, minimum size=0.6cm, inner sep=0] at (0,-5.30) {8};
  \node[anchor=west, font=\scriptsize, text width=7.6cm] at (0.45,-5.30) {Síntesis, cierre actitudinal y bibliografía};
  \fill[iuborange] (8.3,-5.50) rectangle ++(1.53,0.4);
  \node[anchor=west, font=\scriptsize\bfseries] at (9.93,-5.30) {15 min};
  \draw[iubnavy, line width=0.6pt] (8.3,0.45) -- (8.3,-5.70);
\end{tikzpicture}
\end{frame}

% =====================================================================
\bloque{01}{El condensador como elemento que almacena energía en el campo eléctrico y su lugar en los circuitos}{20 min}{iubblue}
% =====================================================================

% ---------------------------------------------------------------------
\begin{frame}{Por qué estudiar el condensador}
\begin{columns}[c,onlytextwidth]
  \column{0.57\textwidth}
\small
\begin{itemize}\setlength{\itemsep}{4pt}
  \item Un capacitor, también llamado condensador, es un dispositivo que guarda carga eléctrica y, con ella, crea un campo eléctrico que a su vez almacena energía $[1]$.
  \item En todos los casos el principio físico es el mismo: dos superficies conductoras separadas por un material aislante, el \textbf{dieléctrico} $[1]$, $[2]$.
  \item Si el circuito externo permanece abierto, el condensador conserva la carga; si se cierra un camino, la energía vuelve al circuito.
\end{itemize}
  \column{0.4\textwidth}
\begin{cardO}{Advertencia de seguridad}
\footnotesize Los capacitores pueden guardar carga durante mucho tiempo después de cortar la alimentación de un circuito $[1]$.
\end{cardO}
\end{columns}
\end{frame}

% ---------------------------------------------------------------------
\begin{frame}{Cadena de funcionamiento del condensador}
\centering
\begin{adjustbox}{max width=\linewidth, max totalheight=6.1cm}
\begin{tikzpicture}
  \node[bG, align=center, font=\scriptsize, text width=1.86cm] (n0) at (1.38,4.63) {Fuente aplica tensión};
  \node[bB, align=center, font=\scriptsize, text width=1.67cm] (n1) at (6.72,4.63) {Placas acumulan carga ±Q};
  \node[bB, align=center, font=\scriptsize, text width=2.16cm] (n2) at (11.88,4.63) {Campo eléctrico en dieléctrico};
  \node[bO, align=center, font=\scriptsize, text width=2.00cm] (n3) at (11.88,1.34) {Energía almacenada};
  \node[dec, align=center, font=\tiny, text width=1.78cm, inner sep=1pt] (n4) at (6.72,1.34) {¿Circuito externo cerrado?};
  \node[bG, align=center, font=\scriptsize, text width=2.52cm] (n5) at (1.38,1.34) {Energía devuelta al circuito};
  \draw[flecha] (n0) -- node[midway, above, fill=white, font=\tiny, inner sep=1pt] {V} (n1);
  \draw[flecha] (n1) -- (n2);
  \draw[flecha] (n2) -- (n3);
  \draw[flecha] (n3) -- (n4);
  \draw[flecha] (n4) -- node[midway, above, fill=white, font=\tiny, inner sep=1pt] {sí} (n5);
  \draw[flecha] (n4) to[bend left=30] node[midway, above, fill=white, font=\tiny, inner sep=1pt] {no: conserva carga} (n3);
\end{tikzpicture}
\end{adjustbox}

\vspace{1mm}{\scriptsize Cadena de funcionamiento del condensador: de la fuente al almacenamiento y la devolución de energía\par}
\end{frame}

% =====================================================================
\bloque{02}{Placas, dieléctrico, C = Q/V y capacitancia de placas paralelas}{35 min}{iuborange}
% =====================================================================

% ---------------------------------------------------------------------
\begin{frame}{{\normalsize Estructura del condensador y definición de capacitancia}}
\begin{columns}[c,onlytextwidth]
  \column{0.57\textwidth}
\small
\begin{itemize}\setlength{\itemsep}{4pt}
  \item Un faradio es la capacitancia de un condensador que guarda un culombio de carga cuando entre sus placas hay un voltio $[1]$.
  \item Con placas de 0.01 m² separadas 1 mil (2.54 × 10\ensuremath{^{-}}\ensuremath{^{5}} m) y mica ($\varepsilon_r = 5$), la capacitancia es $C = (5)(8.85 \times 10^{-12})(0.01) / (2.54 \times 10^{-5}) \approx 0.017\ \mu\text{F}$ $[1]$.
\end{itemize}
  \column{0.4\textwidth}
\begin{caja}{iubblue}{Ecuación clave}
\footnotesize \centering\small \begin{adjustbox}{max width=\linewidth}$\displaystyle C = \frac{Q}{V}$\end{adjustbox}
\end{caja}
\end{columns}
\end{frame}

% ---------------------------------------------------------------------
\begin{frame}{{\normalsize Ecuaciones: Estructura del condensador y definición}}
\begin{columns}[c,onlytextwidth]
  \column{0.55\textwidth}
\begin{cajaplana}{iubblue}
\centering\small \begin{adjustbox}{max width=\linewidth}$\displaystyle \begin{gathered}C = \frac{\varepsilon_r \, \varepsilon_0 \, A}{d} \\ \varepsilon_0 = 8.85 \times 10^{-12}\ \text{F/m}\end{gathered}$\end{adjustbox}
\end{cajaplana}
  \column{0.42\textwidth}
\begin{cardO}{Error frecuente}
\footnotesize Confundir la carga $Q$ (culombios) con la capacitancia $C$ (faradios).
\end{cardO}
\end{columns}
\end{frame}

% ---------------------------------------------------------------------
\begin{frame}{{\normalsize Constantes dieléctricas típicas de materiales usados}}
\centering\scriptsize
\renewcommand{\arraystretch}{1.35}
\rowcolors{2}{iubnavy!6}{white}
\begin{adjustbox}{max width=\linewidth, max totalheight=5.6cm}
\begin{tabularx}{\textwidth}{>{\raggedright\arraybackslash}X>{\raggedright\arraybackslash}X}
  \rowcolor{iubnavy}
  \hd{Material} & \hd{Constante dieléctrica \ensuremath{\epsilon}r típica}\\
  Aire (vacío) & 1.0\\
  Teflón & 2.0\\
  Papel parafinado & 2.5\\
  Aceite & 4.0\\
  Mica & 5.0\\
  Vidrio & 7.5\\
  Cerámica & 1200\\
\end{tabularx}
\end{adjustbox}
\vspace{2mm}
{\scriptsize Constantes dieléctricas típicas de materiales usados en condensadores $[1]$\par}
\end{frame}

% =====================================================================
\bloque{03}{Separación de cargas, rigidez dieléctrica, tensión nominal y energía W = CV\texttwosuperior{}/2}{35 min}{iubgreen}
% =====================================================================

% ---------------------------------------------------------------------
\begin{frame}{Comportamiento interno}
\begin{columns}[c,onlytextwidth]
  \column{0.57\textwidth}
\small
\begin{itemize}\setlength{\itemsep}{4pt}
  \item Lo que ocupa el espacio entre ambas es un \textbf{campo eléctrico}, y en él reside la energía almacenada $[1]$.
  \item Tras la ruptura, el condensador se comporta prácticamente como un conductor.
  \item El dieléctrico real tampoco es un aislante perfecto.
\end{itemize}
  \column{0.4\textwidth}
\begin{caja}{iubblue}{Ecuación clave}
\footnotesize \centering\small \begin{adjustbox}{max width=\linewidth}$\displaystyle \mathcal{E} = \frac{V}{d}$\end{adjustbox}
\end{caja}
\end{columns}
\end{frame}

% ---------------------------------------------------------------------
\begin{frame}{Ecuaciones: Comportamiento interno}
\begin{columns}[c,onlytextwidth]
  \column{0.55\textwidth}
\begin{cajaplana}{iubblue}
\centering\small \begin{adjustbox}{max width=\linewidth}$\displaystyle W = \frac{1}{2} C V^{2}$\end{adjustbox}
\end{cajaplana}
  \column{0.42\textwidth}
\begin{cardO}{Nota pedagógica}
\footnotesize Las 0.147 J del ejemplo parecen poca energía, pero liberadas en microsegundos a través de un destornillador o de los dedos producen chispa y quemaduras.
\end{cardO}
\end{columns}
\end{frame}

% ---------------------------------------------------------------------
\begin{frame}{Energía almacenada en función de la tensión}
\begin{columns}[c,onlytextwidth]
  \column{0.62\textwidth}
\centering
\begin{tikzpicture}
  \begin{axis}[width=8.4cm, height=5.7cm, xmin=0, xmax=30, ymin=0, ymax=0.22419,
      xlabel={Tensión entre terminales V (V)}, ylabel={Energía W (J)},
      label style={font=\scriptsize, text=iubnavy}, tick label style={font=\tiny, text=iubnavy},
      axis line style={iubnavy}, grid=major, grid style={iubnavy!12},
      legend style={font=\tiny, fill=white}, legend pos=north east]
    \addplot[line width=1.4pt, iubblue] coordinates {(0,0) (0.40134,8.0536e-06) (0.80268,3.2214e-05) (1.204,7.2482e-05) (1.6054,0.00012886) (2.0067,0.00020134) (2.408,0.00028993) (2.8094,0.00039463) (3.2107,0.00051543) (3.612,0.00065234) (4.0134,0.00080536) (4.4147,0.00097449) (4.8161,0.0011597) (5.2174,0.0013611) (5.6187,0.0015785) (6.0201,0.0018121) (6.4214,0.0020617) (6.8227,0.0023275) (7.2241,0.0026094) (7.6254,0.0029074) (8.0268,0.0032214) (8.4281,0.0035516) (8.8294,0.0038979) (9.2308,0.0042604) (9.6321,0.0046389) (10.033,0.0050335) (10.435,0.0054442) (10.836,0.0058711) (11.237,0.006314) (11.639,0.0067731) (12.04,0.0072482) (12.441,0.0077395) (12.843,0.0082469) (13.244,0.0087704) (13.645,0.00931) (14.047,0.0098657) (14.448,0.010437) (14.849,0.011025) (15.251,0.011629) (15.652,0.01225) (16.054,0.012886) (16.455,0.013538) (16.856,0.014207) (17.258,0.014891) (17.659,0.015592) (18.06,0.016309) (18.462,0.017041) (18.863,0.01779) (19.264,0.018555) (19.666,0.019337) (20.067,0.020134) (20.468,0.020947) (20.87,0.021777) (21.271,0.022623) (21.672,0.023484) (22.074,0.024362) (22.475,0.025256) (22.876,0.026166) (23.278,0.027092) (23.679,0.028035) (24.08,0.028993) (24.482,0.029967) (24.883,0.030958) (25.284,0.031965) (25.686,0.032988) (26.087,0.034026) (26.488,0.035081) (26.89,0.036153) (27.291,0.03724) (27.692,0.038343) (28.094,0.039463) (28.495,0.040598) (28.896,0.04175) (29.298,0.042918) (29.699,0.044102) (30,0.045)};
    \addlegendentry{C = 100 µF}
    \addplot[line width=1.4pt, iuborange] coordinates {(0,0) (0.40134,3.7852e-05) (0.80268,0.00015141) (1.204,0.00034067) (1.6054,0.00060563) (2.0067,0.0009463) (2.408,0.0013627) (2.8094,0.0018547) (3.2107,0.0024225) (3.612,0.003066) (4.0134,0.0037852) (4.4147,0.0045801) (4.8161,0.0054507) (5.2174,0.006397) (5.6187,0.007419) (6.0201,0.0085167) (6.4214,0.0096901) (6.8227,0.010939) (7.2241,0.012264) (7.6254,0.013665) (8.0268,0.015141) (8.4281,0.016693) (8.8294,0.01832) (9.2308,0.020024) (9.6321,0.021803) (10.033,0.023657) (10.435,0.025588) (10.836,0.027594) (11.237,0.029676) (11.639,0.031833) (12.04,0.034067) (12.441,0.036376) (12.843,0.03876) (13.244,0.041221) (13.645,0.043757) (14.047,0.046369) (14.448,0.049056) (14.849,0.051819) (15.251,0.054658) (15.652,0.057573) (16.054,0.060563) (16.455,0.063629) (16.856,0.066771) (17.258,0.069988) (17.659,0.073281) (18.06,0.07665) (18.462,0.080095) (18.863,0.083615) (19.264,0.087211) (19.666,0.090882) (20.067,0.09463) (20.468,0.098453) (20.87,0.10235) (21.271,0.10633) (21.672,0.11038) (22.074,0.1145) (22.475,0.1187) (22.876,0.12298) (23.278,0.12733) (23.679,0.13176) (24.08,0.13627) (24.482,0.14085) (24.883,0.1455) (25.284,0.15023) (25.686,0.15504) (26.087,0.15992) (26.488,0.16488) (26.89,0.16992) (27.291,0.17503) (27.692,0.18021) (28.094,0.18547) (28.495,0.19081) (28.896,0.19622) (29.298,0.20171) (29.699,0.20728) (30,0.2115)};
    \addlegendentry{C = 470 µF}
    \addplot[line width=1pt, dashed, iubnavy!70] coordinates {(25,0) (25,0.22419)};
    \addlegendentry{Tensión nominal 25 V}
  \end{axis}
\end{tikzpicture}
  \column{0.35\textwidth}
\begin{caja}{iubblue}{Lectura de la gráfica}
\footnotesize La gráfica muestra las dos parábolas: a igual tensión, el condensador de mayor capacitancia almacena proporcionalmente más energía, y en ambos la energía crece con rapidez al acercarse a la tensión nominal.
\end{caja}
\end{columns}
\end{frame}

% =====================================================================
\bloque{04}{Cerámicos, de película, de mica, electrolíticos de aluminio y tantalio, variables}{35 min}{iubblue}
% =====================================================================

% ---------------------------------------------------------------------
\begin{frame}{{\normalsize Tipos de condensadores según su dieléctrico y construcción}}
\begin{columns}[c,onlytextwidth]
  \column{0.57\textwidth}
\small
\begin{itemize}\setlength{\itemsep}{4pt}
  \item Igual que los resistores, los condensadores se agrupan en dos grandes categorías: \textbf{fijos} y \textbf{variables} $[2]$.
  \item Los \textbf{condensadores cerámicos} aprovechan la constante dieléctrica muy alta de la cerámica (1200 es un valor típico) para lograr valores comparativamente altos en tamaños pequeños $[1]$.
  \item La mayoría tienen menos de 1 µF, aunque algunos llegan a 100 µF $[1]$.
\end{itemize}
  \column{0.4\textwidth}
\begin{cardO}{Advertencia de seguridad}
\footnotesize Un condensador electrolítico conectado con la polaridad invertida puede explotar y provocar lesiones $[1]$.
\end{cardO}
\end{columns}
\end{frame}

% ---------------------------------------------------------------------
\begin{frame}{{\normalsize Clasificación de los condensadores según su construcción}}
\centering
\begin{adjustbox}{max width=\linewidth, max totalheight=6.1cm}
\begin{tikzpicture}
  \node[bG, align=center, font=\scriptsize, text width=3.09cm] (n0) at (5.01,5.15) {Condensadores};
  \node[bB, align=center, font=\scriptsize, text width=2.36cm] (n1) at (4.48,3.00) {Fijos no polarizados};
  \node[bB, align=center, font=\scriptsize, text width=3.42cm] (n2) at (8.19,3.00) {Fijos polarizados (electrolíticos)};
  \node[bB, align=center, font=\scriptsize, text width=2.36cm] (n3) at (1.30,3.00) {Variables y trimmers};
  \node[bO, align=center, font=\scriptsize, text width=3.07cm] (n4) at (3.51,0.65) {Mica, cerámica, película};
  \node[bO, align=center, font=\scriptsize, text width=2.18cm] (n5) at (6.96,0.65) {Aluminio y tantalio};
  \draw[flecha] (n0) -- (n1);
  \draw[flecha] (n0) -- (n2);
  \draw[flecha] (n0) -- (n3);
  \draw[flecha] (n1) -- node[midway, above, fill=white, font=\tiny, inner sep=1pt] {según dieléctrico} (n4);
  \draw[flecha] (n2) -- node[midway, above, fill=white, font=\tiny, inner sep=1pt] {según óxido} (n5);
\end{tikzpicture}
\end{adjustbox}

\vspace{1mm}{\scriptsize Clasificación de los condensadores según su construcción y su dieléctrico\par}
\end{frame}

% ---------------------------------------------------------------------
\begin{frame}{Comparación de los tipos de condensadores fijos}
\centering\scriptsize
\renewcommand{\arraystretch}{1.35}
\rowcolors{2}{iubnavy!6}{white}
\begin{adjustbox}{max width=\linewidth, max totalheight=5.6cm}
\begin{tabularx}{\textwidth}{>{\raggedright\arraybackslash}X>{\raggedright\arraybackslash}X>{\raggedright\arraybackslash}X>{\raggedright\arraybackslash}X>{\raggedright\arraybackslash}X}
  \rowcolor{iubnavy}
  \hd{Tipo} & \hd{Dieléctrico} & \hd{Intervalo típico} & \hd{Tensión} & \hd{Polarizado}\\
  Mica & Mica (\ensuremath{\epsilon}r \ensuremath{\approx} 5) & 1 pF a 0.1 µF & 100 V a 2500 V & No\\
  Cerámico & Cerámica (\ensuremath{\epsilon}r \ensuremath{\approx} 1200) & 1 pF a 2.2 µF & hasta 6 kV & No\\
  Película plástica & Poliéster, polipropileno, policarbonato & mayoría < 1 µF & según serie & No\\
  Electrolítico de aluminio & Óxido de aluminio & 1 µF a > 200 000 µF & \ensuremath{\leq} 350 V típico & Sí\\
  Electrolítico de tantalio & Pentóxido de tantalio & valores medios en poco volumen & bajas & Sí\\
\end{tabularx}
\end{adjustbox}
\vspace{2mm}
{\scriptsize Comparación de los tipos de condensadores fijos $[1]$\par}
\end{frame}

% =====================================================================
\bloque{05}{Lectura del valor, la tolerancia y la tensión en el cuerpo del condensador}{30 min}{iuborange}
% =====================================================================

% ---------------------------------------------------------------------
\begin{frame}{{\normalsize Identificación: marcado, código numérico y tolerancia}}
\begin{columns}[c,onlytextwidth]
  \column{0.57\textwidth}
\small
\begin{itemize}\setlength{\itemsep}{4pt}
  \item Identificar un condensador es leer en su cuerpo tres datos: la capacitancia nominal, la tolerancia y la tensión de trabajo.
  \item En los condensadores grandes, como los electrolíticos, el valor se imprime completo: «470 µF 25 V», con el signo menos sobre la franja del terminal negativo.
\end{itemize}
  \column{0.4\textwidth}
\begin{caja}{iubblue}{Ecuación clave}
\footnotesize \centering\small \begin{adjustbox}{max width=\linewidth}$\displaystyle C_{n} = (10\,a + b) \times 10^{\,m}\ \text{pF}$\end{adjustbox}
\end{caja}
\end{columns}
\end{frame}

% ---------------------------------------------------------------------
\begin{frame}{Ecuaciones: Identificación}
\begin{columns}[c,onlytextwidth]
  \column{0.55\textwidth}
\begin{cajaplana}{iubblue}
\centering\small \begin{adjustbox}{max width=\linewidth}$\displaystyle \begin{gathered}C_{\min} = C_{n}\,(1 - t) \\ C_{\max} = C_{n}\,(1 + t)\end{gathered}$\end{adjustbox}
\end{cajaplana}
  \column{0.42\textwidth}
\begin{cardO}{Error frecuente}
\footnotesize Leer «104» como 104 pF.
\end{cardO}
\end{columns}
\end{frame}

% ---------------------------------------------------------------------
\begin{frame}{Lectura de marcados frecuentes de condensadores}
\centering\scriptsize
\renewcommand{\arraystretch}{1.35}
\rowcolors{2}{iubnavy!6}{white}
\begin{adjustbox}{max width=\linewidth, max totalheight=5.6cm}
\begin{tabularx}{\textwidth}{>{\raggedright\arraybackslash}X>{\raggedright\arraybackslash}X>{\raggedright\arraybackslash}X>{\raggedright\arraybackslash}X}
  \rowcolor{iubnavy}
  \hd{Marcado} & \hd{Valor nominal} & \hd{Tolerancia} & \hd{Otros datos}\\
  103 & 10 000 pF = 10 nF & no indicada & cerámico o película\\
  104K 100V & 100 000 pF = 0.1 µF & ±10 \% & 100 V de trabajo\\
  472J & 4700 pF = 4.7 nF & ±5 \% & —\\
  .01 & 0.01 µF = 10 nF & no indicada & unidad implícita µF\\
  330 (cerámico) & 330 pF & no indicada & unidad implícita pF\\
  NP0 & — & — & coeficiente de temperatura nulo\\
  470 µF 25 WVDC & 470 µF & según fabricante & electrolítico, 25 V cd\\
\end{tabularx}
\end{adjustbox}
\vspace{2mm}
{\scriptsize Lectura de marcados frecuentes de condensadores\par}
\end{frame}

% =====================================================================
\bloque{06}{Faradio y sus submúltiplos, conversiones y símbolos normalizados}{25 min}{iubgreen}
% =====================================================================

% ---------------------------------------------------------------------
\begin{frame}{{\normalsize Unidades, múltiplos, submúltiplos, simbología y nomenclatura}}
\begin{columns}[c,onlytextwidth]
  \column{0.57\textwidth}
\small
\begin{itemize}\setlength{\itemsep}{4pt}
  \item Es una unidad enorme para la electrónica y para la mayoría de las aplicaciones de potencia: los valores reales se expresan casi siempre con submúltiplos.
  \item Con el nanofaradio como paso intermedio, cada salto entre unidades consecutivas es un factor 1000.
  \item La \textbf{simbología} del condensador reproduce su estructura: dos trazos paralelos que representan las placas, con un terminal a cada lado $[1]$.
\end{itemize}
  \column{0.4\textwidth}
\begin{cardO}{Error frecuente}
\footnotesize Confundir la «m» minúscula (mili, $10^{-3}$) con la «µ» (micro, $10^{-6}$) o con la «M» mayúscula (mega, $10^{6}$).
\end{cardO}
\end{columns}
\end{frame}

% ---------------------------------------------------------------------
\begin{frame}{Submúltiplos del faradio y reglas de conversión}
\centering\scriptsize
\renewcommand{\arraystretch}{1.35}
\rowcolors{2}{iubnavy!6}{white}
\begin{adjustbox}{max width=\linewidth, max totalheight=5.6cm}
\begin{tabularx}{\textwidth}{>{\raggedright\arraybackslash}X>{\raggedright\arraybackslash}X>{\raggedright\arraybackslash}X>{\raggedright\arraybackslash}X}
  \rowcolor{iubnavy}
  \hd{Unidad} & \hd{Símbolo} & \hd{Valor en F} & \hd{Para pasar a la unidad inmediatamente menor}\\
  milifaradio & mF & 10\ensuremath{^{-}}³ & × 1000 (a µF)\\
  microfaradio & µF & 10\ensuremath{^{-}}\ensuremath{^{6}} & × 1000 (a nF)\\
  nanofaradio & nF & 10\ensuremath{^{-}}\ensuremath{^{9}} & × 1000 (a pF)\\
  picofaradio & pF & 10\ensuremath{^{-}}¹² & —\\
\end{tabularx}
\end{adjustbox}
\vspace{2mm}
{\scriptsize Submúltiplos del faradio y reglas de conversión $[1]$\par}
\end{frame}

% ---------------------------------------------------------------------
\begin{frame}{{\normalsize Símbolos y nomenclatura usados en planos y en el cuerpo}}
\centering\scriptsize
\renewcommand{\arraystretch}{1.35}
\rowcolors{2}{iubnavy!6}{white}
\begin{adjustbox}{max width=\linewidth, max totalheight=5.6cm}
\begin{tabularx}{\textwidth}{>{\raggedright\arraybackslash}X>{\raggedright\arraybackslash}X>{\raggedright\arraybackslash}X}
  \rowcolor{iubnavy}
  \hd{Elemento} & \hd{Representación o marca} & \hd{Significado}\\
  Condensador fijo & dos trazos paralelos & no polarizado\\
  Condensador polarizado & trazo recto y trazo curvo, o signo + & placa recta positiva\\
  Condensador variable & símbolo con flecha diagonal & capacitancia ajustable\\
  Designador & C1, C2, C3… & referencia del componente en el plano\\
  MF, MFD, uF & µF & microfaradio en rotulación abreviada\\
  WV, WVDC & tensión de trabajo & tensión continua admisible\\
\end{tabularx}
\end{adjustbox}
\vspace{2mm}
{\scriptsize Símbolos y nomenclatura usados en planos y en el cuerpo de los condensadores\par}
\end{frame}

% =====================================================================
\bloque{07}{Decodificar el marcado, calcular el intervalo de tolerancia y contrastarlo con la medición}{45 min}{iubblue}
% =====================================================================

% ---------------------------------------------------------------------
\begin{frame}{{\normalsize Identificación y verificación de los condensadores de un kit}}
\begin{columns}[c,onlytextwidth]
  \column{0.57\textwidth}
\small
\begin{itemize}\setlength{\itemsep}{4pt}
  \item Antes de la primera práctica de circuitos capacitivos, el técnico de laboratorio debe dejar listo un kit de cinco condensadores para cada grupo de trabajo.
  \item Para C1, el código 223 da $C_n = 22 \times 10^3\ \text{pF} = 22\ \text{nF}$ y la letra J indica ±5 \%, así que el intervalo es de $22 \times 0.95 = 20.9\ \text{nF}$ a $22 \times 1.05 = 23.1\ \text{nF}$.
  \item Cuatro componentes están dentro de su tolerancia.
\end{itemize}
  \column{0.4\textwidth}
\begin{caja}{iubblue}{Ecuación clave}
\footnotesize \centering\small \begin{adjustbox}{max width=\linewidth}$\displaystyle e_{\%} = \frac{C_{m} - C_{n}}{C_{n}} \times 100$\end{adjustbox}
\end{caja}
\end{columns}
\end{frame}

% ---------------------------------------------------------------------
\begin{frame}{{\normalsize Identificación y verificación de los condensadores del kit}}
\centering\tiny
\renewcommand{\arraystretch}{1.35}
\rowcolors{2}{iubnavy!6}{white}
\begin{adjustbox}{max width=\linewidth, max totalheight=5.6cm}
\begin{tabularx}{\textwidth}{>{\raggedright\arraybackslash}X>{\raggedright\arraybackslash}X>{\raggedright\arraybackslash}X>{\raggedright\arraybackslash}X>{\raggedright\arraybackslash}X>{\raggedright\arraybackslash}X>{\raggedright\arraybackslash}X>{\raggedright\arraybackslash}X>{\raggedright\arraybackslash}X}
  \rowcolor{iubnavy}
  \hd{Ref.} & \hd{Marcado} & \hd{Tipo} & \hd{Cn} & \hd{Tolerancia} & \hd{Intervalo admisible} & \hd{Cm} & \hd{e (\%)} & \hd{Decisión}\\
  C1 & 223J & Cerámico multicapa & 22 nF & ±5 \% & 20.9 a 23.1 nF & 22.4 nF & +1.8 & Acepta\\
  C2 & 104K 100V & Película & 100 nF & ±10 \% & 90 a 110 nF & 97.8 nF & \ensuremath{-}2.2 & Acepta\\
  C3 & 472K & Cerámico de disco & 4.7 nF & ±10 \% & 4.23 a 5.17 nF & 4.52 nF & \ensuremath{-}3.8 & Acepta\\
  C4 & 470 µF 25 V & Electrolítico & 470 µF & ±20 \% & 376 a 564 µF & 512 µF & +8.9 & Acepta\\
  C5 & 100 µF 16 V & Electrolítico reutilizado & 100 µF & ±20 \% & 80 a 120 µF & 71.3 µF & \ensuremath{-}28.7 & Rechaza\\
\end{tabularx}
\end{adjustbox}
\vspace{2mm}
{\scriptsize Identificación y verificación de los condensadores del kit (lecturas de ejemplo)\par}
\end{frame}

% =====================================================================
\bloque{08}{Síntesis, cierre actitudinal y bibliografía}{15 min}{iuborange}
% =====================================================================

% ---------------------------------------------------------------------
\begin{frame}{Síntesis de la sesión}
\begin{columns}[T,onlytextwidth]
  \column{0.32\textwidth}
\begin{cardB}{Estructura del condensador y definición}
\scriptsize Placas, dieléctrico, C = Q/V y capacitancia de placas paralelas
\end{cardB}
  \column{0.32\textwidth}
\begin{cardO}{Comportamiento interno}
\scriptsize Separación de cargas, rigidez dieléctrica, tensión nominal y energía W = CV²/2
\end{cardO}
  \column{0.32\textwidth}
\begin{cardG}{Tipos de condensadores según}
\scriptsize Cerámicos, de película, de mica, electrolíticos de aluminio y tantalio, variables
\end{cardG}
\end{columns}
\vspace{2mm}
\begin{columns}[T,onlytextwidth]
  \column{0.32\textwidth}
\begin{cardB}{Identificación}
\scriptsize Lectura del valor, la tolerancia y la tensión en el cuerpo del condensador
\end{cardB}
  \column{0.32\textwidth}
\begin{cardO}{Unidades, múltiplos, submúltiplos}
\scriptsize Faradio y sus submúltiplos, conversiones y símbolos normalizados
\end{cardO}
  \column{0.32\textwidth}
\begin{cardG}{Identificación y verificación}
\scriptsize Decodificar el marcado, calcular el intervalo de tolerancia y contrastarlo con la medición
\end{cardG}
\end{columns}
\end{frame}

% ---------------------------------------------------------------------
\begin{frame}{Reflexión para llevar}
\centering
\begin{tikzpicture}
  \node[fill=iubnavy, text=white, rounded corners=6pt, text width=11.5cm, inner sep=12pt, align=center, font=\normalsize] (c) at (0,0) {\itshape «Verificar y descargar cada condensador antes de entregarlo al grupo es un acto de responsabilidad con los compañeros: un electrolítico cargado o mal identificado puede lastimar a alguien o arruinar las mediciones de todo el equipo en la práctica siguiente.»};
  \node[fill=iubyellow, text=iubnavy, rounded corners=3pt, font=\scriptsize\bfseries, inner sep=4pt, anchor=south west] at ([yshift=-4pt]c.north west) {Componente actitudinal};
\end{tikzpicture}
\end{frame}

% ---------------------------------------------------------------------
\begin{frame}{Referencias bibliográficas}
\small
\begin{itemize}\setlength{\itemsep}{8pt}
  \item[{$[1]$}] T. L. Floyd, Principios de circuitos eléctricos, 8.\textordfeminine{} ed. México: Pearson Educación, 2007.
  \item[{$[2]$}] R. L. Boylestad, Introducción al análisis de circuitos, 10.\textordfeminine{} ed. México: Pearson Educación, 2004.
\end{itemize}
\end{frame}

% ---------------------------------------------------------------------
%  CIERRE
% ---------------------------------------------------------------------
\begin{frame}[plain]
\begin{tikzpicture}[remember picture, overlay]
  \fill[iubnavy] (current page.north west) rectangle (current page.south east);
  \node[anchor=north west, text=iubyellow, font=\bfseries\fontsize{40}{44}\selectfont]
    at ([xshift=1.2cm,yshift=-1.6cm]current page.north west) {\textexclamdown{}Gracias!};
  \node[anchor=north west, text=white, font=\normalsize, text width=14cm, align=left]
    at ([xshift=1.2cm,yshift=-3.4cm]current page.north west)
    {IEL04 \textperiodcentered{} Circuitos Eléctricos I};
  \node[anchor=north west, fill=iubblue, text=white, rounded corners=1.5mm,
        font=\small\bfseries, inner sep=2.2mm, text width=11.5cm]
    at ([xshift=1.2cm,yshift=-4.4cm]current page.north west)
    {Próxima sesión \textperiodcentered{} Semana 2: Bobinas o inductores: estructura, comportamiento interno, identificación y tipos};
  \node[anchor=south west, text=iubyellow, font=\small, align=left]
    at ([xshift=1.2cm,yshift=0.9cm]current page.south west)
    {Ing. Sergio Martinez Castilla \quad · \quad smartinezc@unibarranquilla.edu.co};
  \fill[iubblue]   ([xshift=1.2cm,yshift=0.55cm]current page.south west) rectangle ++(1.6cm,0.15cm);
  \fill[iuborange] ([xshift=2.9cm,yshift=0.55cm]current page.south west) rectangle ++(1.6cm,0.15cm);
  \fill[iubgreen]  ([xshift=4.6cm,yshift=0.55cm]current page.south west) rectangle ++(1.6cm,0.15cm);
\end{tikzpicture}
\end{frame}

\end{document}
